\chapter{事件、四维量与运动光源}
\section{从测量程序确定事件}\label{th:m1events}
惯性系中，一个事件的位置由静止于该系的尺确定，时间由事件所在地已经同步的钟给出。爱因斯坦同步的约定是：从甲钟发出的光在乙处立即反射，若甲钟发射、接收时刻为$t_1,t_2$，就把乙钟在反射事件的读数校准为$(t_1+t_2)/2$。这一定义用到了同一惯性系内的双程光速$c$。

取$S'$相对$S$沿$x$正向以速度$V$运动，两系原点在事件$O$重合并归零。记
\[
 \beta=V/c,\qquad \gamma=(1-\beta^2)^{-1/2}.
\]
从线性变换$x'=a(x-Vt)$出发，要求$x=\pm ct$在另一系仍有$x'=\pm ct'$，可得$t'=a(t-Vx/c^2)$。逆变换与正变换的形式相同，只把$V$改成$-V$；把两式相乘要求恢复原坐标，得到$a^2(1-\beta^2)=1$。取保持时间方向的正根，得到
\begin{equation}\label{eq:m1lorentz}
 ct'=\gamma(ct-\beta x),\qquad x'=\gamma(x-\beta ct),\qquad y'=y,\quad z'=z.
\end{equation}
逆变换则为$ct=\gamma(ct'+\beta x')$、$x=\gamma(x'+\beta ct')$。这些式子同时适用于任意两事件的坐标差，即使两系坐标原点没有这样选取。

\Diagram{fig:m1coordinates}{同一事件在两惯性系中的坐标；虚线为过$P$的两组等时线}{
\STAxes
\draw[axis](0,0)--(2.5,4.1)node[right]{$ct'$};\draw[axis](0,0)--(4.5,2.7)node[right]{$x'$};
\fill(2.2,3.2)circle(1.3pt)node[above left]{$P$};\draw[guide](0,3.2)--(3.8,3.2);\draw[guide](.3,2.06)--(3.6,4.04);\node[below left]at(0,0){$O$};}
图\ref{fig:m1coordinates}中的$x'$轴是$t'=0$，斜率为$ct/x=\beta$；$ct'$轴是$x'=0$，斜率为$ct/x=1/\beta$。图上两轴的欧氏长度不能直接作为另一系的坐标刻度。若某杆端点的世界线为$x_A(t),x_B(t)$，$S$测得的杆长是$x_B(t)-x_A(t)$，两端事件必须位于同一条水平线上。照片则来自若干不同的发光时刻。

\subsection{同时性、钟慢与尺缩的条件}
在$S$中同时发生的两事件满足$\Delta t=0$，但
\[
 \Delta t'=-\gamma V\Delta x/c^2.
\]
只有$\Delta x=0$或$V=0$时仍然同时。对于固定在$S'$的同一只钟，$\Delta x'=0$，逆变换给$\Delta t=\gamma\Delta t'$，这里$\Delta t'$是这只钟经历的固有时。若测$S'$静止杆的长度$l_0=\Delta x'$，必须另取$\Delta t=0$的一对端点事件；正变换给$l_0=\gamma\Delta x$，所以$l=l_0/\gamma$。钟慢沿同一只钟取事件，尺缩沿同一条等时线取端点；使用公式时应保留各自的事件约束。

\begin{example}\label{ex:m1simultaneous}
在$S$中，左右两个信号灯相距$l$，在$t=0$同时闪光。$S'$沿正向以$0.6c$运动，原点在左灯闪光事件重合。求两闪光事件在$S'$中的时间差、位置差，并说明哪一个先发生。
\Diagram{fig:m1lampsq}{同一水平线上的两个闪光事件}{\STAxes\fill(0,0)circle(1.5pt)node[below]{$A$};\fill(3,0)circle(1.5pt)node[below]{$B$};\draw[guide](0,2)--(3,2);\draw[curve](3,0)--(3,3.6);\node[right]at(3,3.4){右灯};\node[below]at(1.5,0){$l$};}
\end{example}
\begin{solution}
图\ref{fig:m1lampsq}给出$A:(0,0)$、$B:(0,l)$，坐标次序取$(ct,x)$。代入式\eqref{eq:m1lorentz}，$A'=(0,0)$，而$B'$为
\[
 ct'_B=-\gamma\beta l=-\frac34l,\qquad x'_B=\gamma l=\frac54l.
\]
故$\Delta t'=t'_B-t'_A=-3l/(4c)$，右灯先闪。这里$5l/4$是两次闪光事件的空间坐标差，不是$S'$测得的两灯间距。
\Diagram{fig:m1lampsa}{在$S'$中取同时事件才能测量灯距，坐标轴按相同比例作图}{\draw[axis](-.3,0)--(4.3,0)node[right]{$x'$};\draw[axis](0,-2.8)--(0,2.2)node[above]{$ct'$};\draw[curve](1.5,-2.5)--(-.9,1.5);\draw[curve](3.9,-2.5)--(1.5,1.5);\fill(0,0)circle(1.4pt)node[left]{$A$};\fill(3.75,-2.25)circle(1.4pt)node[right]{$B$};\draw[guide](0,0)--(2.4,0);\fill(2.4,0)circle(1.4pt)node[above right]{$C$};}
图\ref{fig:m1lampsa}中的$C$取在右灯世界线上，并满足$t'_C=0$。右灯在$S$中$x=l$，由$t'=0$得$t=Vl/c^2$，再代入$x'=\gamma(l-Vt)$，得$x'_C=l/\gamma=4l/5$。这才是$S'$的灯距。
\end{solution}

\begin{example}\label{ex:m1rods}
杆$AB$静止于$S$，杆$A'B'$静止于$S'$，两杆固有长均为$l=ct_0$，$S'$沿正向以$V$运动。$A$与$A'$相遇于$M$，这两只端点钟归零；$B$与$B'$相遇于$N$，那两只钟归零。所有钟此后正常走时。$M$处从$A'$向$B'$发出光，光在$B'$立即反射，向$A'$返回。求光出射时经过$B$、在$B'$反射、返回时经过$B$、到达$A'$及到达$A$这五个事件中，各相应端点钟的读数，并讨论返回光线遇到各端点的次序。
\Diagram{fig:m1rodsq}{两杆及局部归零事件，图示为$S$中的一个同时截面}{\draw[axis](-.4,0)--(6,0)node[right]{$x$};\draw[curve](0,.7)--(3,.7);\draw[curve](.8,1.7)--(3.2,1.7);\foreach\x/\y/\t in{0/.7/A,3/.7/B,.8/1.7/A',3.2/1.7/B'}{\fill(\x,\y)circle(1.3pt)node[above]{$\t$};}\draw[axis](3.7,1.7)--(5,1.7)node[right]{$V$};\draw[<->](0,.3)--(3,.3)node[midway,below]{$l$};\draw[<->](.8,2.2)--(3.2,2.2)node[midway,above]{$l/\gamma$};}
\end{example}
\begin{solution}
以$M$为两坐标系原点。图\ref{fig:m1rodsq}中$B$固定在$x=l$，$B'$的世界线为$x=Vt+l/\gamma$。相遇条件给
\[
 t_N=\frac{l}{V}(1-\gamma^{-1})=\frac{t_0}{\beta}(1-\gamma^{-1}).
\]
因此$B$钟读$t-t_N$，$A$钟读$t$。在$S'$中，$A'$钟读$t'$，$B'$钟读$t'-t'_N$，其中
\begin{align*}
 t'_N&=\gamma(t_N-\beta t_0)
 =\frac{\gamma t_0}{\beta}(1-\gamma^{-1}-\beta^2)\\
 &=\frac{t_0}{\beta}(\gamma^{-1}-1)=-t_N.
\end{align*}
最后一步用$1-\beta^2=\gamma^{-2}$。两个不同的局部归零操作没有把各杆上的两只钟同步。
\Diagram{fig:m1rodsa}{光的全部交点；$P,Q$在$B$世界线上，$R$为反射事件}{
\draw[axis](-.2,0)--(4.9,0)node[right]{$x/l$};\draw[axis](0,-.15)--(0,6.4)node[above]{$ct/l$};
\draw[curve](1,0)--(1,6);\node[above]at(1,6){$B$};\draw[curve](0,0)--(3.6,6);\node[right]at(3.6,6){$A'$};\draw[curve](.8,0)--(4.4,6);\node[right]at(4.4,6){$B'$};
\draw[very thick,->](0,0)--(2,2);\draw[very thick,->](2,2)--(0,4);
\foreach\x/\y/\t in{0/0/M,1/.333/N,1/1/P,2/2/R,1/3/Q,1.5/2.5/T,0/4/U}{\fill(\x,\y)circle(1.4pt)node[above left]{$\t$};}}
图\ref{fig:m1rodsa}按$\beta=0.6$画出比例，计算保留一般$\beta$。出射光$x=ct$，故$P$为$t_P=t_0$。光在$S'$中走过固定长度$l$后到$R$，所以$(ct'_R,x'_R)=(l,l)$，逆变换给
\[
 t_R=\gamma(1+\beta)t_0,\quad x_R=ct_R.
\]
反射光在$S$中的方程是$x=2ct_R-ct$。与$x=l$相交得$t_Q=2t_R-t_0$；与$x=0$相交得$t_U=2t_R$。在$S'$中，光往返的总路程为$2l$，故回到$A'$的事件$T$满足$t'_T=2t_0$。于是所求各钟读数为
\begin{align*}
 B(P):&\quad t_0-t_N,\\
 B'(R):&\quad t_0-t'_N=t_0+t_N,\\
 B(Q):&\quad 2\gamma(1+\beta)t_0-t_0-t_N,\\
 A'(T):&\quad 2t_0,\\
 A(U):&\quad 2\gamma(1+\beta)t_0.
\end{align*}
返回途中，$t_T=2\gamma t_0$，故$t_Q-t_T=(2\gamma\beta-1)t_0$。令此差为零，得$4\beta^2=1-\beta^2$，即$\beta=1/\sqrt5$。因此$0<\beta<1/\sqrt5$时依次遇到$B,A',A$；$\beta>1/\sqrt5$时依次遇到$A',B,A$；临界时$B$与$A'$恰在光路上相遇。若取$\beta=0.6$，则$\gamma=1.25$、$t_N=t_0/3$，五个钟面读数依次为$2t_0/3,4t_0/3,8t_0/3,2t_0,4t_0$，而$T$发生在$Q$之前。事件排序使用同一参考系的坐标时间，不能直接比较未同步的端点钟。
\end{solution}

\section{四维矢量与固有量}\label{th:m1four}
以$x^\mu=(ct,x,y,z)$表示四维位置，取度规$g_{\mu\nu}=\operatorname{diag}(1,-1,-1,-1)$。降指标得到$x_\mu=(ct,-x,-y,-z)$，重复上下指标表示求和。两个四维矢量的内积为
\[
 A\cdot B=A^\mu B_\mu=A^0B^0-A^1B^1-A^2B^2-A^3B^3.
\]
将式\eqref{eq:m1lorentz}代入可逐项验证$\Delta s^2=c^2\Delta t^2-\Delta\vct r^2$不变。矩阵语言是$\Lambda^{\mathsf T}g\Lambda=g$。四维矢量按同一个$\Lambda$变换，其内积因此不变；仅发现某个量的平方不变，还不足以证明任意构造的四个数具有正确变换律。

\subsection{类时、类空与类光间隔}
若$\Delta s^2>0$且$\Delta t>0$，两事件可由亚光速质点依次到达，并存在使两事件发生在同一位置的惯性系。这不要求参与两个事件的不同物体都在该系静止。若$\Delta s^2=0$且事件间隔非零，两事件可由真空光线连接。若$\Delta s^2<0$，存在使两事件同时的惯性系，部分惯性系会交换其时间顺序；它们不能由不超过光速的信号建立因果关系。

对有质量质点的世界线，微小间隔满足
\begin{equation}\label{eq:m1proper}
 c^2\dd\tau^2=c^2\dd t^2-\dd\vct r^2,
 \qquad \dd\tau=\dd t\sqrt{1-v^2/c^2}.
\end{equation}
弯曲世界线的固有时要沿轨迹积分，不能把端点间隔的平方根当作整段固有时。两点间匀速直线世界线的固有时最大：在两端点共同的静止系中，$\int\sqrt{1-v^2/c^2}\dd t\leq\int\dd t$，等号只在全程$v=0$时成立。

\begin{example}\label{ex:m1interval}
两事件在$S$中的坐标差为$c\Delta t=3L$、$\Delta x=5L$。求使它们同时的惯性系及该系中的距离，说明是否有使它们同地的惯性系。
\Diagram{fig:m1intervalq}{类空分离的两个事件位于光锥外}{\STAxes\draw[guide](0,0)--(4.2,4.2);\fill(4,2.4)circle(1.4pt)node[right]{$B$};\node[below left]at(0,0){$A$};\draw[curve](0,0)--(4,2.4);}
\end{example}
\begin{solution}
图\ref{fig:m1intervalq}中两事件的间隔为
\[
 \Delta s^2=(3L)^2-(5L)^2=-16L^2<0.
\]
这是类空间隔，允许改变事件的时间先后。沿$x$作速度$V$的变换，令两事件同时，条件为
\[
 0=c\Delta t'=\gamma(3L-5\beta L),\qquad \beta=V/c.
\]
因此$\beta=3/5$、$\gamma=1/\sqrt{1-9/25}=5/4$。在该系中空间距离为
\[
 \Delta x'=\frac54\left(5-\frac35\times3\right)L=4L.
\]
\Diagram{fig:m1intervala}{同时系中的两事件，间隔绝对值对应空间距离}{\draw[axis](-.3,0)--(5,0)node[right]{$x'$};\draw[axis](0,-.2)--(0,2.4)node[above]{$ct'$};\fill(0,0)circle(1.4pt)node[below]{$A$};\fill(4,0)circle(1.4pt)node[below]{$B$};\draw[<->](0,.6)--(4,.6)node[midway,above]{$4L$};}
图\ref{fig:m1intervala}的$x'$轴通过两事件；它代表$t'=0$的同时线。此时不变量为$-(\Delta x')^2=-16L^2$，与原系相同。若要求两事件同地，则须令
\[
 0=\Delta x'=\gamma(5L-3\beta L),\qquad \beta=5/3.
\]
这个速度不属于惯性系之间的亚光速变换。也可直接从不变量排除：同地时应有$\Delta s^2=c^2\Delta t'^2\geq0$，与题设矛盾。使两事件同时的$x$向速度是唯一的；若还允许横向变换，可在已经同时的系中继续沿垂直于事件分离方向变换，仍保持同时。
\end{solution}

\subsection{体积和密度}
静止小体积$V_0$在运动方向缩短，实验室的同时截面给$V=V_0/\gamma$。若物质粒子的静质量总和为$M_0$，实验室单位体积所含静质量为$\rho_m=\gamma\rho_{m0}$；它与实验室能量密度不是同一个量。无压力冷物质的能量密度为$u=\gamma^2\rho_{m0}c^2$，因为每个粒子的能量还多一因子$\gamma$。有压力或内部应力的介质须使用能量动量张量，不能继续套用这个冷尘结果。

\begin{example}\label{ex:m1density}
一束粒子在共同静止系的数密度为$n_0$，每粒子静质量$m$、电荷$q$。粒子束在实验室沿$x$正向以$v$运动。求数密度、静质量密度、电荷密度、电流密度及能量密度。
\Diagram{fig:m1densityq}{粒子束的静止体积和实验室同时截面}{\draw[curve](0,0)rectangle(2.7,1.4);\draw[curve](4,0)rectangle(6,1.4);\foreach\x in{.4,1.2,2}{\fill(\x,.7)circle(2pt);}\foreach\x in{4.3,4.9,5.5}{\fill(\x,.7)circle(2pt);}\node[above]at(1.35,1.4){$V_0$};\node[above]at(5,1.4){$V_0/\gamma$};\draw[axis](4,-.5)--(5.8,-.5)node[right]{$v$};}
\end{example}
\begin{solution}
按图\ref{fig:m1densityq}，同一批$N=n_0V_0$粒子占据$V_0/\gamma$，所以$n=\gamma n_0$。每粒子的静质量、电荷不因变换改变，故
\[
 \rho_m=\gamma n_0m,\quad \rho_e=\gamma n_0q,\quad j_x=\rho_ev.
\]
实验室每粒子能量$\gamma mc^2$，于是$u=n\gamma mc^2=\gamma^2n_0mc^2$。
\Diagram{fig:m1densitya}{穿过截面的粒子数决定电流，能量通量另乘单粒子能量}{\draw[curve](3,-.3)--(3,1.8);\node[above]at(3,1.8){面积$A$};\draw[curve](0,0)rectangle(3,1.4);\draw[axis](.5,.7)--(2.5,.7)node[midway,above]{$v$};\node[below]at(1.5,0){$v\dd t$};\node[right]at(3.3,.6){$\dd Q=qnAv\dd t$};}
图\ref{fig:m1densitya}中时间$\dd t$内穿过面积$A$的粒子来自长度$v\dd t$的小柱，故$\dd Q/A\dd t=qnv$。同理能流为$uv$。密度和通量必须使用同一参考系的体积、时间与单粒子物理量。
\end{solution}

\section{速度、光行差与方向变换}\label{th:m1velocity}
对式\eqref{eq:m1lorentz}求微分，并除以$\dd t'$，得
\begin{equation}\label{eq:m1vel}
 v'_x=\frac{v_x-V}{1-Vv_x/c^2},\qquad
 v'_{y,z}=\frac{v_{y,z}}{\gamma(1-Vv_x/c^2)}.
\end{equation}
分母来自不同时性的时间变换，不能只对位移做变换。若$v<c$，四速度的时间分量为正，因而分母始终正。由直接代数还可得
\[
 1-\frac{v'^2}{c^2}=\frac{1-v^2/c^2}{\gamma^2(1-Vv_x/c^2)^2}>0,
\]
所以亚光速速度经洛伦兹变换仍为亚光速。

真空光线在$S'$中的方向角为$\theta'$，相对$x'$正向测量。将$v'_x=c\cos\theta'$、$v'_y=c\sin\theta'$代入逆速度变换，得到
\begin{equation}\label{eq:m1aberration}
 \cos\theta=\frac{\cos\theta'+\beta}{1+\beta\cos\theta'},\quad
 \sin\theta=\frac{\sin\theta'}{\gamma(1+\beta\cos\theta')}.
\end{equation}
两式保留了角度象限，优于单独使用正切式。还可写成半角形式
\[
 \tan\frac\theta2=\sqrt{\frac{1-\beta}{1+\beta}}\tan\frac{\theta'}2.
\]
前半球压缩是方向变换的结果，并不意味着在光源静止系中发光机制变得不对称。

\begin{example}\label{ex:m1aberration}
在光源静止系，光子沿与运动方向垂直的方向发出。求实验室中的光线方向；再求静止系前半球在实验室对应的圆锥及其立体角比例。
\Diagram{fig:m1aberrationq}{源静止系的前半球及边缘光线}{\draw[axis](-1.8,0)--(2.6,0)node[right]{$x'$};\draw[axis](0,-1.6)--(0,2)node[above]{$y'$};\draw[curve](0,-1.4)arc(-90:90:1.4);\draw[very thick,->](0,0)--(0,1.5);\draw[axis](2,.8)--(3.2,.8)node[right]{$V$};}
\end{example}
\begin{solution}
取源系光子四动量为$(E'/c,E'\cos\theta'/c,E'\sin\theta'/c,0)$。逆变换给
\[
 E=\gamma E'(1+\beta\cos\theta'),\quad
 cp_x=\gamma E'(\cos\theta'+\beta),\quad cp_y=E'\sin\theta'.
\]
由光子$p=E/c$，图\ref{fig:m1aberrationq}中光线的方向满足
\[
 \cos\theta=\frac{\cos\theta'+\beta}{1+\beta\cos\theta'},\quad
 \sin\theta=\frac{\sin\theta'}{\gamma(1+\beta\cos\theta')}.
\]
令$\theta'=\pi/2$，边缘方向便为$\cos\theta_c=\beta$、$\sin\theta_c=1/\gamma$。
\Diagram{fig:m1aberrationa}{实验室中的前向光锥，边缘满足$\cos\theta_c=\beta$}{\draw[axis](-.5,0)--(4,0)node[right]{$x$};\draw[curve](0,0)--(3,1.5);\draw[curve](0,0)--(3,-1.5);\draw[guide](3,1.5)arc(90:270:.25 and 1.5);\draw[->](.9,0)arc(0:26.565:.9);\node[right]at(.9,.45){$\theta_c$};\draw[very thick,->](0,0)--(2.4,1.2);}
设$\mu'=\cos\theta'$，对方向变换求导，得到
\[
 \frac{\dd\cos\theta}{\dd\mu'}
 =\frac{(1+\beta\mu')-\beta(\mu'+\beta)}{(1+\beta\mu')^2}
 =\frac{1-\beta^2}{(1+\beta\mu')^2}>0.
\]
源系前半球$0\leq\mu'\leq1$因此完整映为图\ref{fig:m1aberrationa}中的$\beta\leq\cos\theta\leq1$。绕$x$轴旋转此角域，其立体角为
\[
 \Omega_c=2\pi\int_0^{\theta_c}\sin\theta\dd\theta
 =2\pi[-\cos\theta]_0^{\theta_c}
 =2\pi(1-\beta),\qquad \frac{\Omega_c}{4\pi}=\frac{1-\beta}{2}.
\]
高速时$1-\beta\simeq1/(2\gamma^2)$，故$\Omega_c\simeq\pi/\gamma^2$、$\theta_c\simeq1/\gamma$。角域缩小描述的是方向变换；同一角域中的光子数和能量份额还须分别按发射分布与频移加权。
\end{solution}

\begin{example}\label{ex:m1tilt}
一根静止在$S'$的细杆，其端点坐标差为$(l_0\cos\alpha',l_0\sin\alpha')$。求$S$测得的杆长和倾角，并与一束沿杆方向传播的光的方向变换比较。
\Diagram{fig:m1tiltq}{杆的方向由两个同时端点确定}{\draw[axis](-.3,0)--(4,0)node[right]{$x'$};\draw[axis](0,-.2)--(0,2.8)node[above]{$y'$};\draw[very thick](0,0)--(3,2);\node[above]at(1.5,1.1){$l_0$};\draw(.8,0)arc(0:33.69:.8);\node at(1.2,.35){$\alpha'$};}
\end{example}
\begin{solution}
图\ref{fig:m1tiltq}中的杆静止，所以其横向分量不变、纵向分量缩短。对$S$中同时的端点事件，$\Delta x'=\gamma\Delta x$，得到
\[
 l_x=\frac{l_0\cos\alpha'}\gamma,\quad l_y=l_0\sin\alpha',\quad
 l=l_0\sqrt{\frac{\cos^2\alpha'}{\gamma^2}+\sin^2\alpha'}.
\]
故$\tan\alpha=\gamma\tan\alpha'$，实际象限由两个分量确定。
\Diagram{fig:m1tilta}{杆的等时方向与沿杆光线的传播方向有不同定义}{\draw[axis](-.3,0)--(4.6,0)node[right]{$x$};\draw[axis](0,-.2)--(0,3)node[above]{$y$};\draw[very thick](0,0)--(2,2);\node[above left]at(2,2){杆};\draw[->,thick](0,0)--(3.7,1.3);\node[right]at(3.7,1.3){光};\draw[guide](2,0)--(2,2);}
光的两个事件是同一光子先后经过两处的事件，须用式\eqref{eq:m1aberration}，并取$\theta'=\alpha'$，得到$\tan\theta=\sin\alpha'/[\gamma(\cos\alpha'+\beta)]$。图\ref{fig:m1tilta}展示了这种差别。发射后杆仍沿$x$方向运动，因此另一系中的光路方向一般不同于杆的同时截面方向。
\end{solution}

\section{四速度、四动量与四力}\label{th:m1dynamics}
固有时是不变量，因此位置四矢量对固有时的导数仍是四矢量：
\[
 u^\mu=\frac{\dd x^\mu}{\dd\tau}=\gamma_v(c,\vct v),
 \quad u\cdot u=c^2.
\]
静质量不变的质点有$p^\mu=mu^\mu=(E/c,\vct p)$，即$E=\gamma_vmc^2$、$\vct p=\gamma_vm\vct v$。它们满足
\begin{equation}\label{eq:m1shell}
 E^2-c^2p^2=m^2c^4.
\end{equation}
静能不是人为增加在动能上的常数：能量与动量须共同按四矢量变换，静能由它们的变换关系决定。低速展开$\gamma_v=1+v^2/(2c^2)+\cdots$后，才恢复$E=mc^2+mv^2/2+\cdots$。

令三力$\vct f=\dd\vct p/\dd t$。对式\eqref{eq:m1shell}求导，$E\dd E=c^2\vct p\cdot\dd\vct p$，并利用$c^2\vct p/E=\vct v$，得
\[
 \frac{\dd E}{\dd t}=\vct f\cdot\vct v=\mathcal P.
\]
因而四力为$F^\mu=\dd p^\mu/\dd\tau=\gamma_v(\mathcal P/c,\vct f)$。四加速度$a^\mu=\dd u^\mu/\dd\tau$与四速度正交：对$u^2=c^2$求导，得$u\cdot a=0$。静质量固定时$F^\mu=ma^\mu$，从而$u\cdot F=0$。若研究吸收能量导致静质量变化的物体，则$F^\mu=m a^\mu+u^\mu\dd m/\dd\tau$，不能略去后一项。

\begin{example}\label{ex:m1restforce}
质点瞬时沿$x$方向运动，实验室三力为$(f_\parallel,f_\perp,0)$。求四力的平方，并求瞬时静止系的三力。
\Diagram{fig:m1restforceq}{速度和三力的平行、垂直分解}{\draw[axis](0,0)--(3.7,0)node[right]{$\vct v$};\draw[axis](0,0)--(2.4,1.8)node[right]{$\vct f$};\draw[guide](2.4,0)--(2.4,1.8)--(0,1.8);\node[below]at(1.2,0){$f_\parallel$};\node[left]at(0,.9){$f_\perp$};}
\end{example}
\begin{solution}
四力定义为$F^\mu=\dd p^\mu/\dd\tau$，而三力为$\vct f=\dd\vct p/\dd t$。图\ref{fig:m1restforceq}中速度沿$x$，功率为$\mathcal P=\vct f\cdot\vct v=f_\parallel v$，故
\[
 F^\mu=\gamma_v(f_\parallel v/c,f_\parallel,f_\perp,0).
\]
取度规$(+,-,-,-)$，逐项求平方：
\begin{align*}
 F^2&=\gamma_v^2[f_\parallel^2(\beta_v^2-1)-f_\perp^2]\\
 &=-f_\parallel^2-\gamma_v^2f_\perp^2.
\end{align*}
这里使用了$\gamma_v^2(1-\beta_v^2)=1$。
\Diagram{fig:m1restforcea}{瞬时静止系中的力；此时四力时间分量为零}{\draw[axis](0,0)--(3.6,0)node[right]{$x'$};\draw[axis](0,0)--(0,3)node[above]{$y'$};\draw[very thick,->](0,0)--(2,2.4)node[right]{$\vct f_0$};\draw[guide](2,0)--(2,2.4)--(0,2.4);\node[below]at(1,0){$f_\parallel$};\node[left]at(0,1.2){$\gamma_vf_\perp$};}
沿$x$以速度$v$变换至图\ref{fig:m1restforcea}的瞬时静止系。时间、纵向和横向分量分别为
\begin{align*}
 F'^0&=\gamma_v(F^0-\beta_vF^x)=0,\\
 F'^x&=\gamma_v(F^x-\beta_vF^0)
 =\gamma_v^2f_\parallel(1-\beta_v^2)=f_\parallel,\\
 F'^y&=F^y=\gamma_v f_\perp.
\end{align*}
该事件处$v'=0$、$\gamma_{v'}=1$，所以四力空间分量就是三力，得到$f_{0\parallel}=f_\parallel$、$f_{0\perp}=\gamma_vf_\perp$，且$f_0^2=-F^2$。恒静质量粒子的四力因而是类空或零矢量。时间分量为零表示该系中瞬时功率为零；粒子静止是充分条件，力与速度垂直时也满足这一条件。
\end{solution}

\section{运动光源的角分布}
\subsection{频率与立体角的雅可比}
光子的四动量为$(h\nu/c,(h\nu/c)\vct n)$。对时间分量作变换，得
\[
 \nu=\gamma\nu'(1+\beta\cos\theta')
 =\frac{\nu'}{\gamma(1-\beta\cos\theta)}=\delta\nu',
\]
其中$\delta$称多普勒因子。例\ref{ex:m1aberration}已给出光行差商式的求导过程；以$\mu'=\cos\theta'$表示，结果为
\[
 \dd\mu=\frac{1-\beta^2}{(1+\beta\mu')^2}\dd\mu',
 \qquad \dd\Omega=\delta^{-2}\dd\Omega'.
\]
方位角不变，因而立体角的变换只需处理这一导数。

\begin{example}\label{ex:m1power}
点光源在自身静止系中各向同性发光，单位固有时、单位立体角的辐射能为$I_0$，总功率$P_0=4\pi I_0$。光源以速度$V$沿$x$运动。求按实验室发射时间定义的角功率$I(\theta)$及总功率，并与远处探测器按接收时间定义的角功率比较。
\Diagram{fig:m1powerq}{同一小束光在两系中占据不同立体角}{\draw[axis](-.5,0)--(4.2,0)node[right]{$V$};\draw[curve](0,0)--(3,1);\draw[curve](0,0)--(3,1.7);\draw[curve](3,1)arc(-90:90:.18 and .35);\node[right]at(3.3,1.4){$\dd\Omega$};\draw[guide](0,0)--(3,1.35);\node at(1.2,.24){$\theta$};}
\end{example}
\begin{solution}
先在图\ref{fig:m1powerq}的一束光中数同一批光子，光子数不变，单光子能量乘$\delta$，故$\dd E=\delta\dd E'$。光源两个发射事件位于同一世界线上，实验室发射时间$\dd t=\gamma\dd\tau$。因此
\begin{equation}\label{eq:m1angular}
 I(\theta)=\frac{\dd E}{\dd t\dd\Omega}
 =I_0\frac\delta\gamma\frac{\dd\Omega'}{\dd\Omega}
 =\frac{I_0}{\gamma^4(1-\beta\cos\theta)^3}.
\end{equation}
积分总功率时使用实验室立体角，令$\mu=\cos\theta$：
\begin{align*}
 P&=\frac{2\pi I_0}{\gamma^4}\int_{-1}^{1}\frac{\dd\mu}{(1-\beta\mu)^3}\\
 &=\frac{2\pi I_0}{\gamma^4}
 \left[\frac{1}{2\beta(1-\beta\mu)^2}\right]_{-1}^{1}\\
 &=\frac{2\pi I_0}{\gamma^4}\frac1{2\beta}
 \frac{(1+\beta)^2-(1-\beta)^2}{(1-\beta^2)^2}
 =4\pi I_0.
\end{align*}
原函数由$z=1-\beta\mu$、$\dd\mu=-\dd z/\beta$及$\int z^{-3}\dd z=-z^{-2}/2$得到。以上写法取$\beta\ne0$；静止情形直接积分$I_0$，也得$P=4\pi I_0$。
这个结果还有整体四动量解释：静止系一次各向同性发光的总动量为零，实验室总能量乘$\gamma$，发射时间也乘$\gamma$。
\Diagram{fig:m1powera}{两次发射到达固定远处探测器的路程差}{\draw[axis](-.3,0)--(4,0)node[right]{$x$};\fill(0,0)circle(1.4pt)node[below]{$S_1$};\fill(1,0)circle(1.4pt)node[below]{$S_2$};\fill(4,2)circle(1.4pt)node[right]{$D$};\draw[curve](0,0)--(4,2);\draw[curve](1,0)--(4,2);\draw[guide](1,0)--(.8,.4);\node[above]at(2,1.5){$\theta$近似不变};}
图\ref{fig:m1powera}中两个发射点相距$V\dd t$，第二束路程缩短$V\dd t\cos\theta$，因此
\[
 \dd t_{\rm arr}=(1-\beta\cos\theta)\dd t,
 \quad I_{\rm arr}(\theta)=I_0\delta^4.
\]
这里探测器足够远，方向在两次发射间近似不变。接收时间的角功率不能再对同一个接收时刻的全天球方向简单积分而解释成同一批发射光的总功率；不同方向对应不同的发射时段。
\end{solution}

\section{坐标测量中的事件选择}
\subsection{从光锥得到变换系数}
惯性系的时空均匀性要求：同一个坐标差采用同一变换，与起点位置、起始时刻无关。坐标关系因而具有仿射形式；将原点对齐后，变换为线性。先把原点对齐，沿$x$的变换可写$x'=a x+bct$、$ct'=d x+ect$。$S'$原点在$S$中满足$x=\beta ct$，因此$b=-a\beta$。

分别代入向右、向左的两条原点光线。$x=ct$要求$x'=ct'$，所以$a+b=d+e$；$x=-ct$要求$x'=-ct'$，所以$-a+b=d-e$。两式相加、相减给$d=b=-a\beta$、$e=a$。至此光锥不变只确定了整体比例$a$，尚未确定其数值。用相对性原理要求逆变换在交换两系后只有$\beta$变号，正逆矩阵相乘便给$a^2(1-\beta^2)=1$。连续到$\beta=0$的恒等变换，选$a=\gamma>0$。

这一推导也说明为什么单独要求光速不变还不能任意省略比例系数。若同时允许改变坐标单位，光锥形状不变但标尺标度也会改变；本书要求两系采用相同物理单位和互易的标准钟尺。

横向坐标不变可由对称性理解：沿$x$的运动没有选出横向某一方向。设横向统一缩放为$k(\beta)$，正逆变换要求$k(\beta)k(-\beta)=1$，空间反射对称给两者相等，连续正根为1。这排除了把横向尺寸也按$\gamma$缩短的做法。

\subsection{同一事件、同一物体与同一时刻}
坐标题有三种常见约束，各自消去不同未知量。两条世界线相遇要求同一事件的$t,x$都相同；跟踪同一物体是在其世界线上取两事件；同时测长度则是在某一系固定$t$后取不同物体的事件。只有把题设约束先转为这些关系，代入变换才有明确意义。

固定一个钟的两个事件满足$x'=\text{常数}$，因而$\Delta x=V\Delta t$。代入时间变换得$\Delta t'=\gamma\Delta t(1-\beta^2)=\Delta t/\gamma$。若改成在$S$同时读取两只钟，便有$\Delta t=0$而不是$\Delta x=V\Delta t$；相同的式子转而给$\Delta t'=-\gamma V\Delta x/c^2$。两种情境都由同一变换式产生，不能把它们当作两个可随意替换的“钟慢规则”。

\Diagram{fig:m1pairings}{沿同一钟的两事件与沿同一等时线的两事件，是不同配对}{\draw[axis](-.2,0)--(4.4,0)node[right]{$x$};\draw[axis](0,0)--(0,4.5)node[above]{$ct$};\draw[curve](0,0)--(2.4,4);\draw[curve](1,0)--(3.4,4);\fill(.6,1)circle(1.5pt)node[left]{$A$};\fill(1.8,3)circle(1.5pt)node[left]{$B$};\fill(2.8,3)circle(1.5pt)node[right]{$C$};\draw[guide](1.8,3)--(2.8,3);}
图\ref{fig:m1pairings}中$A,B$在同一钟世界线上，$B,C$在同一实验室等时线上。$AB$适合求钟经历的时间，$BC$适合求实验室钟距。若问题还给定钟面被设过零，最后须在坐标时间上加钟的设置常数，不能让这个常数反过来改变世界线。

\subsection{四钟模型的完整事件坐标}
例\ref{ex:m1rods}可把所有坐标集中写成无量纲表。定义$g=\gamma(1+\beta)$、$a=(1-1/\gamma)/\beta$，坐标单位分别为$t_0,l$。$A$钟读$t/t_0$，$B$钟读$t/t_0-a$；$A'$钟读$t'/t_0$，$B'$钟读$t'/t_0+a$。表中每行都代表一个事件，列之间只做坐标变换。
\begin{center}\small
\begin{tabular}{ccccc}\toprule
事件&$t/t_0$&$x/l$&$t'/t_0$&$x'/l$\\\midrule
$M$&0&0&0&0\\
$N$&$a$&1&$-a$&1\\
$P$&1&1&$1/g$&$1/g$\\
$R$&$g$&$g$&1&1\\
$Q$&$2g-1$&1&$\gamma(2g-1-\beta)$&$\gamma[1-\beta(2g-1)]$\\
$T$&$2\gamma$&$2\gamma\beta$&2&0\\
$U$&$2g$&0&$2\gamma g$&$-2\gamma\beta g$\\\bottomrule
\end{tabular}
\end{center}
$P$行使用$\gamma(1-\beta)=1/g$；$R$行则由源系最简单的$(1,1)$变来。$N$行的$t'_N=-t_N$是本题两杆等固有长和特定归零事件共同造成的结果，换成不等长两杆后一般不再成立。

表中的$Q,T$两行也重现例\ref{ex:m1rods}已求出的返回次序及临界速度$\beta=1/\sqrt5$。几何上比较$x_T=2\gamma\beta l$与固定$B$位置$l$即可判断：反射光向左走，先遇到更右侧的对象。$\beta=.6$时$A'$在右，故先$T$后$Q$。

这些次序针对有因果联系的返回光线，同一系或另一系都会给出相同先后。局部归零后的两只不同钟却可能读出反向大小，因此“表上较大的读数”不能单独代替事件次序。$\beta\to0$时$a\simeq\beta/2\to0$，公式连续趋于静止往返，但正好$\beta=0$时两杆端点一直重合，所谓第二次相遇归零事件需要另作规定。

\section{四维量的几何和角分布的测度}
\subsection{任意方向的纯推动}
沿单位矢量$\hat{\vct V}$分解空间矢量$\vct r=\vct r_\parallel+\vct r_\perp$。对平行分量使用一维变换，垂直分量保持不变，得到
\[
 ct'=\gamma\left(ct-\frac{\vct V\cdot\vct r}{c}\right),\quad
 \vct r'=\vct r+\left[\frac{\gamma-1}{V^2}(\vct V\cdot\vct r)-\gamma t\right]\vct V.
\]
$V=0$的极限取恒等映射。对任意四矢量，只需把$ct$换成其时间分量、把$\vct r$换成空间分量。四动量的能量变换因而是$E'=\gamma(E-\vct V\cdot\vct p)$。这里$\vct V$是坐标系速度，$\vct p$是被研究粒子的动量，二者没有理由总是平行。

洛伦兹标量在任意允许的惯性变换下都保持数值，如固有时、静质量和两个四矢量的缩并。某一分量在一次特殊变换中恰好不变，还不足以判定它是标量；光子能量就可能在特定方向和推动速度下不变，换一种变换后又改变。

类时四矢量的空间分量可在某个系全为零。非零类光四矢量的闵可夫斯基平方为零，时间与空间分量都不能通过有限惯性变换消去；真空光子因而没有静止系。类空四矢量可以把时间分量变为零，却不能把空间分量全消去。这些分类取决于闵可夫斯基内积；时空图的纸面长度不表示四矢量模。

\subsection{单位角、单位立体角与环带}
轴对称分布中，单位立体角强度$I(\theta)$与极角环带内功率不同。环带固体角$\dd\Omega=2\pi\sin\theta\dd\theta$，故$\dd P/\dd\theta=2\pi I(\theta)\sin\theta$。即使$I(\theta)$在$\theta=0$最大，有限环带功率的密度也因$\sin\theta$在轴上为零而不在那里最大。画角分布前须说明纵轴是哪种密度。

\Diagram{fig:m1solidring}{极角环带宽度相同，立体角仍随$\sin\theta$改变}{\draw[curve](0,0)circle(1.7);\draw[axis](0,0)--(2.7,0)node[right]{轴};\draw[curve](0,0)--(1.472,.85);\draw[curve](0,0)--(1.302,1.093);\draw[thick](1.472,.85)arc(30:40:1.7);\node[above]at(.4,2){$\dd\Omega=2\pi\sin\theta\dd\theta$};}
图\ref{fig:m1solidring}以截面短弧表示球面环带。极坐标曲线的半径表示所选的角密度，积分时仍须乘立体角雅可比。运动源的功率份额可先选定一个系中的角域，再变换边界计算，往往能省去完整角分布的积分。

更一般地，源系辐射不必各向同性。设单位固有时发出的总能量为$P_0$，总$x$向动量率为$\dot p'_x$。整个发射小过程的四动量变换给$\dd E=\gamma(\dd E'+\beta c\dd p'_x)$，而发射时间$\dd t=\gamma\dd\tau$，于是实验室总发射功率为
\[
 P=P_0+\beta c\dot p'_x.
\]
各向同性源有$\dot p'_x=0$，所以总功率不变；向后单束光有$\dot p'_x=-P_0/c$，所以$P=P_0(1-\beta)$。两种结果都由同一四动量式得到。总功率不变需要源系净辐射动量为零，各向同性发射满足这一条件。

\subsection{密度变换中的超平面与通量}
粒子数守恒并不表示不同参考系中的数密度相等。计数所用的体积必须处在该系的同时面上；同一批粒子在两个同时面上被截取的位置不同。对共同四速度为$U^\mu$的稀疏粒子流，定义数流$N^\mu=n_0U^\mu$。在实验室中，
\[
 N^\mu=(cn,\vct j_N)=n_0\gamma_v(c,\vct v),\qquad
 n=\gamma_vn_0,\quad \vct j_N=n\vct v.
\]
穿过时空小面元的粒子数为$\dd\mathcal N=N^\mu\dd\Sigma_\mu/c$。实验室等时小体积有$\dd\Sigma_\mu=(\dd V,0,0,0)$，给出$n\dd V$；跨越固定空间面的面元则给出粒子通量。密度与流量是同一四维对象的不同分量，变换时不能仅替换长度而保持时间率不变。

无热运动的粒子束还具有能量动量流$T^{\mu\nu}=n_0mU^\mu U^\nu$。其分量为
\[
 T^{00}=n_0m\gamma_v^2c^2=u,\quad
 cT^{0x}=uv_x,\quad T^{xx}=n_0m\gamma_v^2v_x^2.
\]
最后一项是单位面积、单位时间跨面的$x$动量；即使束在自身系没有压力，运动系也测得定向动量通量。对存在随机热运动的气体，须对每个粒子的速度分布求和，不能再把整个系统视为只有一个共同$U^\mu$的冷束。
\Diagram{fig:m1currentcut}{同一粒子流由不同同时面截取；斜线不代表粒子轨迹的偏折}{\draw[axis](0,0)--(4.5,0)node[right]{$x$};\draw[axis](0,0)--(0,3.7)node[above]{$ct$};\foreach\a in{.4,1.2,2,2.8}{\draw[curve](\a,0)--({\a+.7},3.2);}\draw[guide](.1,1.1)--(3.8,1.1)node[right]{$t=\text{常量}$};\draw[densely dotted](.1,1.6)--(3.8,2.7)node[right]{$t'=\text{常量}$};}
图\ref{fig:m1currentcut}的两条截线各自连接一组不同事件。密度变换应同时指明计数的同时面和跨面通量，两者由同一四流联系。静质量密度$n m$、总能量密度$u$和$u/c^2$分别包含不同的洛伦兹因子，不能互作同义量。

\section{有限装置中的同时事件与雷达测距}
相对论题中的“同时打开”“一起读数”必须绑定某一组静止同步钟。装置的连接条件还可能受信号传播限制。几何上可任意选择一条等时线与若干世界线的交点；物理上是否能用同一触发信号控制这些交点，要另看它们的因果关系。下列模型把这两个层次分开。

\begin{example}\label{ex:m1tunnel}
隧道静止长度$L$，列车固有长也是$L$，列车以$\beta c$进入。地面系在列车中点经过隧道中点时，同时短暂关闭两门并立即打开。证明地面看来列车能在两门之间，并求车系两门关闭的时间差和位置差，解释车系隧道更短却无矛盾。
\Diagram{fig:m1tunnelq}{地面同时关门，列车等时长度为$L/\gamma$}{\draw[axis](-.5,0)--(6,0)node[right]{$x$};\draw[very thick](0,-.8)--(0,1.3);\draw[very thick](5,-.8)--(5,1.3);\draw[curve](.6,.2)rectangle(4.4,.8);\draw[axis](2.3,1.6)--(3.7,1.6)node[right]{$V$};\node[below]at(2.5,-.8){两门距离$L$};}
\end{example}
\begin{solution}
取关门中点事件为$(ct,x)=(0,0)$，两门关闭为$D_L=(0,-L/2)$、$D_R=(0,L/2)$。图\ref{fig:m1tunnelq}的列车两端在该时刻位于$\pm L/(2\gamma)$，严格处于两门之间。变换两关闭事件，得
\[
 ct'_L=\frac{\gamma\beta L}2,\quad x'_L=-\frac{\gamma L}2;
 \qquad ct'_R=-\frac{\gamma\beta L}2,\quad x'_R=\frac{\gamma L}2.
\]
所以右门先关，时间差$t'_R-t'_L=-\gamma\beta L/c$。两次关门的事件位置差是$\gamma L$，并非车系同时隧道长度$L/\gamma$。
\Diagram{fig:m1tunnela}{车系中两关闭事件不同时；虚线为车系等时线}{\draw[axis](-2.6,0)--(3,0)node[right]{$x'$};\draw[axis](0,-2)--(0,2.4)node[above]{$ct'$};\draw[curve](-2,1.2)--(0,-2.1);\draw[curve](0,2.1)--(2,-1.2);\fill(-2,1.2)circle(1.4pt)node[left]{$D_L$};\fill(2,-1.2)circle(1.4pt)node[right]{$D_R$};\draw[guide](-2,1.2)--(1.5,1.2);\draw[guide](-1.5,-1.2)--(2,-1.2);}
图\ref{fig:m1tunnela}中前门关闭时，列车尾部尚在隧道外；后门稍后关闭时，前门早已打开，车头已可通过。车系从未声称两门同时夹住整列车。若要求门持续关闭至列车撞门并停止，则题目变为受力和信号传播问题，不能继续把列车视为始终匀速的刚体。
\end{solution}

\begin{example}\label{ex:m1radar}
观察者固定在$x=0$，匀速飞行的反射镜满足$x=L+Vt$，$L>0,V>0$。观察者在$t=0$发光，光在镜面反射后返回。求反射、返回时刻，以及观察者按雷达约定赋给反射事件的时间和距离。
\Diagram{fig:m1radarq}{出射光追赶向外运动的镜面}{\STAxes\draw[curve](1,0)--(3,4);\node[right]at(3,4){镜};\draw[very thick,->](0,0)--(2,2);\node[below]at(1,0){$L$};}
\end{example}
\begin{solution}
图\ref{fig:m1radarq}的交点满足$ct_R=L+Vt_R$，所以$t_R=L/(c-V)$、$x_R=cL/(c-V)$。反射光以$c$返回原点，返回时刻$t_A=t_R+x_R/c=2L/(c-V)$。
\Diagram{fig:m1radara}{反射事件的雷达时间是发射、返回时刻的平均}{\draw[axis](-.2,0)--(4.6,0)node[right]{$x$};\draw[axis](0,-.2)--(0,5)node[above]{$ct$};\draw[very thick,->](0,0)--(2,2);\draw[very thick,->](2,2)--(0,4);\draw[guide](0,2)--(2,2);\node[left]at(0,2){$ct_{\rm rad}$};\node[right]at(2,2){$R$};\node[left]at(0,4){$ct_A$};}
雷达定义$t_{\rm rad}=(0+t_A)/2=L/(c-V)$，$r_{\rm rad}=ct_A/2=cL/(c-V)$，与图\ref{fig:m1radara}中反射事件的惯性坐标$t_R,x_R$完全一致。雷达同步把距离赋给反射事件，因此这里应与反射时刻的镜距比较。发射时的$L$和接收时的$L+Vt_A$属于另外两个事件。
\end{solution}

\begin{example}\label{ex:m1meeting}
两质点在实验室$t=0$位于$x=0,L$，速度分别为$+u,-w$，其中$0<u,w<c$。求相遇事件及在第一质点静止系中的相遇时间，说明为什么该系的初始两点距离不能直接取$L/\gamma_u$再任意相除。
\Diagram{fig:m1meetingq}{两条匀速世界线的交点}{\draw[axis](-.3,0)--(4.7,0)node[right]{$x$};\draw[axis](0,-.2)--(0,4.5)node[above]{$ct$};\draw[curve](0,0)--(2,3.5);\draw[curve](4,0)--(2,3.5);\node[below]at(4,0){$L$};\node[right]at(2,3.5){$C$};}
\end{example}
\begin{solution}
图\ref{fig:m1meetingq}满足$ut=L-wt$，所以$t_C=L/(u+w)$、$x_C=uL/(u+w)$。第一质点系以$u$运动，且把其初始事件归零，因此
\[
 t'_C=\gamma_u(t_C-ux_C/c^2)=\frac{L}{\gamma_u(u+w)}.
\]
这是第一质点从出发到相遇的固有时。若要在其系中用距离除以速度，必须重新求$t'=0$与第二质点世界线的交点。
\Diagram{fig:m1meetinga}{在第一质点系中，第二质点初始位置取在$t'=0$线上}{\draw[axis](-.2,0)--(4.3,0)node[right]{$x'$};\draw[axis](0,-.2)--(0,4.4)node[above]{$ct'$};\draw[curve](3.4,0)--(0,3.7);\draw[very thick](0,0)--(0,3.7);\node[below]at(3.4,0){$D$};\node[left]at(0,3.7){$C$};}
令第二质点$x=L-wt$满足$t-ux/c^2=0$，得到$t_D=uL/(c^2+uw)$，$x_D=Lc^2/(c^2+uw)$，进而$x'_D=L/[\gamma_u(1+uw/c^2)]$。第二质点相对速度为$-(u+w)/(1+uw/c^2)$。图\ref{fig:m1meetinga}距离除以相对速大小，恰得上式$t'_C$。$L$是两运动物体间的某系等时距离，不是一个共同静止杆的固有长。
\end{solution}

\section{光锥、波面和速度空间}
\begin{example}\label{ex:m1flash}
点光源在事件$O$瞬时闪光。证明实验室固定时刻$t$的光前沿是以$O$空间位置为中心的球面，并说明与运动光源角分布的关系。
\Diagram{fig:m1flashq}{光锥的水平截面是球面，不随之后的光源位置移动}{\draw[axis](-2.4,0)--(2.6,0)node[right]{$x$};\draw[axis](0,-.2)--(0,4)node[above]{$ct$};\draw[curve](0,0)--(2,3.2);\draw[curve](0,0)--(-2,3.2);\draw[curve](0,3.2)ellipse(2 and .4);\draw[guide](0,0)--(.9,3.2);\node[right]at(.9,2.6){源世界线};}
\end{example}
\begin{solution}
每个光前沿事件与$O$的间隔为零。在源静止系$\Delta s'^2=0$，变换不改间隔，故$c^2t^2-x^2-y^2-z^2=0$，取未来支$t>0$便得半径$ct$球面。图\ref{fig:m1flashq}中的源在闪光后已移动，其当前位置不是这一次闪光球面的中心。
\Diagram{fig:m1flasha}{球面形状与能量分布分属几何和光子统计两个层面}{\draw[curve](0,0)circle(1.6);\foreach\a in{0,20,40,60,90,140,210,270,300,330,345}{\draw[->](0,0)--({1.6*cos(\a)},{1.6*sin(\a)});}\draw[axis](0,0)--(2.8,0)node[right]{$V$};\node[below]at(0,-2){球面上能量可不均匀};}
图\ref{fig:m1flasha}用较密箭头示意前向辐射更集中。球面由传播速度和发射事件决定，角能量密度由光行差、单光子频移决定；各向同性波速并不要求各向同性光子分布。对持续运动源，每一次发射有各自的球面，所有球面叠加的稳态场又是另一问题。
\end{solution}

\begin{example}\label{ex:m1lightellipse}
在源系，一次闪光到光子发射后固定源系时刻$t'_0$的全部前沿点组成圆。将这一组事件变到实验室，求它们在$x,y$坐标投影中的曲线，并说明为什么不违背上一例的球面结论。
\Diagram{fig:m1lightellipseq}{本题固定源系时间，所选事件不在实验室同时面上}{\draw[axis](-2.2,0)--(2.5,0)node[right]{$x'$};\draw[axis](0,-2)--(0,2.3)node[above]{$y'$};\draw[curve](0,0)circle(1.5);\node[above]at(1.5,1.5){$r'=ct'_0$};}
\end{example}
\begin{solution}
按图\ref{fig:m1lightellipseq}，令$r_0=ct'_0$，各事件$x'=r_0\cos\theta'$、$y'=r_0\sin\theta'$。逆变换给
\[
 x=\gamma r_0(\cos\theta'+\beta),\quad y=r_0\sin\theta',\quad
 ct=\gamma r_0(1+\beta\cos\theta').
\]
\Diagram{fig:m1lightellipsea}{异时事件的空间投影是平移的椭圆}{\draw[axis](-1.5,0)--(3.8,0)node[right]{$x$};\draw[axis](0,-2)--(0,2.2)node[above]{$y$};\draw[curve](1.322876,0)ellipse(2 and 1.5);\draw[guide](1.322876,-1.5)--(1.322876,1.5);\node[above]at(1.322876,1.6){中心$\gamma\beta r_0$};}
消去$\theta'$，图\ref{fig:m1lightellipsea}的投影满足$(x-\gamma\beta r_0)^2/(\gamma^2r_0^2)+y^2/r_0^2=1$。但这些点的实验室时间随角度变化，范围为$\gamma t'_0(1\pm\beta)$。上一例固定的是同一实验室时间，两题选择的事件集合不同，因此投影为椭圆并不表示实验室真空光速各向异性。
\end{solution}

\begin{example}\label{ex:m1relativeangle}
在$S'$中，两粒子速度大小同为$u$，方向分别与$x'$正向成$\pm\alpha$。求实验室中的夹角，并讨论当$u=c$和$u\ll c$时有何差别。
\Diagram{fig:m1relativeangleq}{关于运动轴对称的两条发射速度}{\draw[axis](0,0)--(3.7,0)node[right]{$V$};\draw[axis](0,0)--(2.5,1.6)node[right]{$u$};\draw[axis](0,0)--(2.5,-1.6)node[right]{$u$};\node at(1,.35){$\alpha$};}
\end{example}
\begin{solution}
图\ref{fig:m1relativeangleq}两粒子都有$u'_x=u\cos\alpha$，横向分量为$\pm u\sin\alpha$。逆速度变换给
\[
 v_x=\frac{u\cos\alpha+V}{1+Vu\cos\alpha/c^2},\quad
 v_y=\frac{\pm u\sin\alpha}{\gamma(1+Vu\cos\alpha/c^2)}.
\]
当变换后的速度非零时，每条速度与$x$的有向角可用$\vartheta=\operatorname{atan2}(u\sin\alpha,\gamma(u\cos\alpha+V))$表示，两线夹角取$2\vartheta$的较小几何夹角。若$u\sin\alpha=0$且$u\cos\alpha+V=0$，粒子在实验室静止，方向及夹角均无定义，不能调用$\operatorname{atan2}(0,0)$。
\Diagram{fig:m1relativeanglea}{同向变换压缩前向发射角，但须保留纵向分量的符号}{\draw[axis](0,0)--(4,0)node[right]{$x$};\draw[axis](0,0)--(3.2,.8);\draw[axis](0,0)--(3.2,-.8);\draw[guide](0,0)--(2,1.5);\draw[guide](0,0)--(2,-1.5);\node[right]at(3.2,.8){实验室};\node[right]at(2,1.5){源系方向示意};}
若$u=c$，此式化为光行差，所有粒子的速率仍为$c$。若$u,V\ll c$，$\gamma\simeq1$且分母近1，恢复伽利略速度矢量平移。图\ref{fig:m1relativeanglea}取$\cos\alpha>0$；若发射近后方使$u\cos\alpha+V<0$，两粒子仍可能向后运动，单独使用正切会误把它们画到前半平面。
\end{solution}

\section{角分布积分与四维微分关系}
\begin{example}\label{ex:m1hemipower}
各向同性静止光源以速度$V$运动，按实验室发射时间定义总功率。求实验室前半球$0\leq\theta\leq\pi/2$收到的功率份额，并与源系前半球光子集合的份额区分。
\Diagram{fig:m1hemipowerq}{实验室前半球是按实验室方向筛选的集合}{\draw[axis](-1.8,0)--(2.6,0)node[right]{$x$};\draw[curve](0,-1.6)arc(-90:90:1.6);\draw[guide](0,-1.8)--(0,1.8);\draw[axis](0,0)--(1.4,.8);}
\end{example}
\begin{solution}
使用式\eqref{eq:m1angular}，图\ref{fig:m1hemipowerq}的前半球积分为
\[
 P_F=\frac{2\pi I_0}{\gamma^4}\int_0^1\frac{\dd\mu}{(1-\beta\mu)^3}
 =\frac{2\pi I_0}{\gamma^4}
 \left[\frac{(1-\beta\mu)^{-2}}{2\beta}\right]_0^1
 =\frac{\pi I_0}{\beta\gamma^4}[(1-\beta)^{-2}-1].
\]
除以总功率$4\pi I_0$，再用$\gamma^{-4}=(1-\beta)^2(1+\beta)^2$，得到
\[
 \frac{P_F}P=\frac{(2-\beta)(1+\beta)^2}{4}
 =\frac{2+3\beta-\beta^3}{4}.
\]
\Diagram{fig:m1hemipowera}{前半球功率份额从一半连续趋向全部}{\draw[axis](0,0)--(4.2,0)node[right]{$\beta$};\draw[axis](0,0)--(0,2.8)node[above]{$P_F/P$};\draw[curve,domain=0:1,samples=70]plot({3.5*\x},{2.3*(2+3*\x-\x^3)/4});\node[left]at(0,1.15){$1/2$};\node[left]at(0,2.3){$1$};\node[below]at(3.5,0){$1$};}
图\ref{fig:m1hemipowera}在$\beta\to0$给$1/2$，在$\beta\to1$趋1。这里边界是实验室$\cos\theta=0$，对应源系$\cos\theta'=-\beta$，包含源系后半球的一部分；不能替换成例\ref{ex:m1coneenergy}中源系前半球所占的$1/2+\beta/4$能量份额。
\end{solution}

\begin{example}\label{ex:m1curvedproper}
粒子在实验室作半径$R$、速率$v$的匀速圆周运动，转一周后回到同一位置。求粒子固有时、四加速度及其不变量，说明速度方向变化是否影响钟慢因子。
\Diagram{fig:m1curvedq}{匀速圆周速度切向、加速度径向}{\draw[curve](0,0)circle(1.5);\fill(1.5,0)circle(2pt);\draw[axis](1.5,0)--(1.5,1.2)node[right]{$v$};\draw[axis](1.5,0)--(.3,0)node[above]{$v^2/R$};}
\end{example}
\begin{solution}
图\ref{fig:m1curvedq}中实验室周期$T=2\pi R/v$。速率恒定，故$\tau=\int_0^T\sqrt{1-v^2/c^2}\dd t=[t/\gamma_v]_0^T=T/\gamma_v$。方向变化不改变局部钟慢因子，固有时仍沿粒子的实际世界线累积。
\Diagram{fig:m1curveda}{圆周世界线在时空中成螺旋，固有时沿世界线累积}{\draw[axis](0,0)--(0,4.2)node[above]{$ct$};\draw[curve,domain=0:720,samples=150]plot({1.5*cos(\x)},{.005*\x+.3*sin(\x)});\node[right]at(1.7,3.6){空间圆周投影示意};}
先对$\gamma_v=(1-v^2/c^2)^{-1/2}$用链式法则求导：$\dot\gamma_v=\gamma_v^3\vct v\cdot\vct a/c^2=0$。再对$u^\mu=(\gamma_vc,\gamma_v\vct v)$逐项求导，
\[
 a^\mu=\gamma_v\frac{\dd u^\mu}{\dd t}
 =\gamma_v(c\dot\gamma_v,\dot\gamma_v\vct v+\gamma_v\vct a)
 =(0,\gamma_v^2\vct a).
\]
所以$a^2=-\gamma_v^4v^4/R^2$，固有加速度$\alpha=\gamma_v^2v^2/R$。图\ref{fig:m1curveda}是螺旋的二维投影，投影交叉点未必是同一时空事件。留在起点的实验室钟走过$T$，粒子钟少走$T(1-1/\gamma_v)$。
\end{solution}

\section{相对速率的标量表达与辐射的局部能量}
\subsection{任意方向的相对速率}
两粒子的实验室速度分别为$\vct u,\vct v$，四速度为$U=\gamma_u(c,\vct u)$、$V=\gamma_v(c,\vct v)$。在第二粒子静止系，$V'=(c,0,0,0)$，故$U'\cdot V'=\gamma_{\rm rel}c^2$。内积不变立即给
\[
 \gamma_{\rm rel}=\gamma_u\gamma_v
 \left(1-\frac{\vct u\cdot\vct v}{c^2}\right),\qquad
 w=c\sqrt{1-\gamma_{\rm rel}^{-2}}.
\]
这里$w$指任一粒子静止系中测得的另一粒子速率，交换两粒子后数值不变。在第三个参考系中作速度矢量差，得到的是该系等时位置差的变化率；它与上述相对速率采用了不同的测量程序。

若两速度共线同向，公式化为$w=|u-v|/(1-uv/c^2)$；共线反向则分子变为$u+v$、分母为$1+uv/c^2$。非共线时，先算夹角余弦形成的点积通常比逐个变换三个分量更直接。若还需要相对运动方向，才补做速度变换；一个标量速率不足以确定方向。

\begin{example}正交运动的两粒子
实验室中$A$沿$x$以$0.8c$运动，$B$沿$y$以$0.6c$运动，两者在原点相遇。求它们的相对速率，并在$B$的静止系中用分量变换核对。
\ExpansionDiagram{fig:dep-m1-relative-q}{相遇事件处的两条实验室速度；虚线表示坐标速度差}{
\draw[->](0,0)--(4,0)node[right]{$x$};\draw[->](0,0)--(0,3)node[above]{$y$};\draw[->,thick](0,0)--(3.2,0)node[midway,below]{$0.8c$};\draw[->,thick](0,0)--(0,2.4)node[midway,left]{$0.6c$};\draw[dashed,->](0,2.4)--(3.2,0);\node[below left]at(0,0){$O$};}
\end{example}
\begin{solution}
图\ref{fig:dep-m1-relative-q}中两速度点积为零，$\gamma_A=5/3$、$\gamma_B=5/4$，所以
\[
 \gamma_{\rm rel}=\frac{25}{12},\qquad
 w=\frac{\sqrt{481}}{25}c\simeq0.8773c.
\]
变到沿$+y$以$0.6c$运动的系，$A$原来的$u_y=0$，速度变换的公共分母为$1-u_yV/c^2=1$。纵向分量变为$u_y'=-0.6c$，横向为$u_x'=u_x/\gamma_B=0.64c$。
\ExpansionDiagram{fig:dep-m1-relative-a}{在$B$静止系中，$A$的横向分量也发生变化}{
\draw[->](0,0)--(3.8,0)node[right]{$x'$};\draw[->](0,-2.8)--(0,1)node[above]{$y'$};\draw[->,thick](0,0)--(2.56,-2.4)node[right]{$\vct u'_A$};\draw[dashed](2.56,0)--(2.56,-2.4)--(0,-2.4);\node[above]at(1.3,0){$0.64c$};\node[left]at(0,-1.4){$-0.6c$};}
由图\ref{fig:dep-m1-relative-a}，$w^2/c^2=0.64^2+0.6^2=481/625$，与内积法一致。若直接使用原图两速度差，会得到$c$；错误来自把实验室的同时性和时间间隔原封不动地带入$B$系。
\end{solution}

\subsection{局部能量密度为什么还与压力有关}
光子束的粒子数四流与单个光子的四动量分别变换。设某源系方向余弦为$\mu$的一小束光在源系能量密度为$\dd u_*$；源系相对实验室沿$+x$运动，速度为$\beta c$。光子能量乘$\gamma(1+\beta\mu)$，同一方向粒子束的数密度也乘这个因子，故
\[
 \dd u=\gamma^2(1+\beta\mu)^2\dd u_*.
\]
这里仍以源系方向为积分标签，尚未把积分测度换为实验室立体角。用同一标签积分可避免重复引入角度雅可比。该束在实验室的$x$方向余弦为$(\mu+\beta)/(1+\beta\mu)$，于是其能流为$c\dd u$乘这个余弦。

对源系各向同性辐射，角平均的权重是$\dd\mu/2$。所需角度矩为
\[
 \langle\mu\rangle=\frac12\int_{-1}^1\mu\dd\mu
 =\left[\frac{\mu^2}{4}\right]_{-1}^1=0,
 \qquad
 \langle\mu^2\rangle=\frac12\int_{-1}^1\mu^2\dd\mu
 =\left[\frac{\mu^3}{6}\right]_{-1}^1=\frac13.
\]
将能量因子平方展开为$1+2\beta\mu+\beta^2\mu^2$，能流的角度因子展开为$(1+\beta\mu)(\mu+\beta)$，逐项取上述平均，得到
\[
 u=\gamma^2u_*(1+\beta^2/3),\qquad
 S_x=\frac43\gamma^2\beta c u_*.
\]
再以方向余弦平方求法向动量通量，有$p_x=\gamma^2u_*(1/3+\beta^2)$；横向压力仍为$u_*/3$。辐射运动后不再各向同性，不能继续令所有方向压力都等于实验室$u/3$。

\begin{example}移动辐射腔中的能量与能流
长方形腔在自身静止系装有各向同性平衡辐射，能量密度$u_*$、体积$\mathcal V_*$。腔相对实验室沿$x$以$0.6c$运动。求辐射在实验室的局部能量密度、能流、纵向压力和等时截面总能量。解释为何最后一个量不等于单独把$u_*\mathcal V_*$乘$\gamma$。
\ExpansionDiagram{fig:dep-m1-radiation-q}{腔系辐射各向同性；实验室中的等时截面另有方向权重}{
\draw(0,0)rectangle(4,2);\foreach\a in{0,45,90,135,180,225,270,315}{\draw[->](2,1)--++(\a:.7);}\draw[->,thick](4.4,1)--(6,1)node[right]{$V$};\node[below]at(2,0){腔系体积$\mathcal V_*$};}
\end{example}
\begin{solution}
图\ref{fig:dep-m1-radiation-q}中的源系角平均给$\gamma=5/4$。代入局部公式，
\[
 u=\frac74u_*,\qquad S_x=\frac54cu_*,\qquad
 p_x=\frac{13}{12}u_*.
\]
横向压力为$u_*/3$。能流虽非零，但$S_x/(cu)=5/7<1$；它是不同方向光子贡献的合成，不是说所有光子共同以某个亚光速运动。
\ExpansionDiagram{fig:dep-m1-radiation-a}{压力和能流共同参与运动辐射的局部描述}{
\draw(0,0)rectangle(3.2,1.8);\draw[->,thick](1,.9)--(4.8,.9)node[right]{$S_x$};\draw[->](3.2,1.4)--(4.2,1.4)node[right]{$p_x$};\draw[->](1.6,1.8)--(1.6,2.6)node[above]{$u_*/3$};\node[below]at(1.6,0){$u=7u_*/4$};}
实验室等时体积为$\mathcal V_*/\gamma=4\mathcal V_*/5$，所以辐射总能量$E_{\rm rad}=7u_*\mathcal V_*/5$，而$\gamma u_*\mathcal V_*=5u_*\mathcal V_*/4$。

两者不同并不违反四动量变换。图\ref{fig:dep-m1-radiation-a}中的腔壁持续承受辐射压力，辐射子系统通过壁面交换动量，并非孤立封闭系统。将腔壁的能量、动量和应力一并纳入，封闭整体的总四动量才按四矢量变换。不能把一团受约束辐射的局部密度积分，无条件当作自由粒子的质量能公式。
\end{solution}

\section{习题与解答}
\begin{exercise}\label{ex:m1anisotropicsource}
光源静止系单位固有时、单位立体角辐射能为$I'(\theta')=I_0(1+a\cos\theta')$，$|a|\leq1$。源沿对称轴以$V$运动。求实验室发射时间角功率和总功率，讨论何时总功率大于静止系值。
\Diagram{fig:m1anisotropicq}{源系辐射含沿运动轴的一阶不对称项}{\draw[axis](-2,0)--(3.1,0)node[right]{$x',V$};\draw[curve,domain=0:360,samples=100]plot({1.2*(1+.6*cos(\x))*cos(\x)},{1.2*(1+.6*cos(\x))*sin(\x)});\node[above]at(0,1.8){图取$a>0$};}
\end{exercise}
\begin{solution}
图\ref{fig:m1anisotropicq}总静止功率仍为$P_0=4\pi I_0$，因为$\cos\theta'$全空间积分为零。对同一批光子使用能量、时间与立体角三因子，
\[
 I(\theta)=\frac{\delta^3}{\gamma}I_0
 \left[1+a\frac{\cos\theta-\beta}{1-\beta\cos\theta}\right].
\]
要积分总功率，改在源系按$\mu'=\cos\theta'$积分更直接。每束的实验室能量乘$\gamma(1+\beta\mu')$，除发射时间$\gamma\dd\tau$，所以
\begin{align*}
 P&=2\pi I_0\int_{-1}^1(1+a\mu')(1+\beta\mu')\dd\mu'\\
 &=2\pi I_0\int_{-1}^1[1+(a+\beta)\mu'+a\beta\mu'^2]\dd\mu'\\
 &=2\pi I_0\left[\mu'+\frac{a+\beta}{2}\mu'^2+
 \frac{a\beta}{3}\mu'^3\right]_{-1}^{1}
 =4\pi I_0(1+a\beta/3).
\end{align*}
\Diagram{fig:m1anisotropica}{总功率的额外项来自源系净辐射动量}{\draw[curve](0,0)rectangle(2.7,1);\node at(1.35,.5){$P_0$};\draw[->](2.9,.5)--(3.7,.5);\node[right]at(3.9,.5){$P=P_0+\beta c\dot p'_x$};\node[below]at(2,-.7){$\dot p'_x=aP_0/(3c)$};}
图\ref{fig:m1anisotropica}中，总功率的修正项由源系净辐射动量决定。$a\beta>0$时实验室功率更大，反号时更小，$a=0$才恢复各向同性源的总功率不变。给定发射机制若长期维持不对称辐射，光源会有反冲；本题假定外部装置维持其匀速运动，或只研究一个瞬时发射小过程。
\end{solution}

\begin{exercise}[选做]\label{ex:m1antitelephone}
假设存在在任意发射者静止系均可沿任意方向以同一速率$u>c$发送的信号。地球$A$固定在$x=0$，飞船$B$满足$x=L+Vt$。$A$在$t=0$发信号，$B$收到后立即在自身系以$-u$回复。求回复到达地球时刻，讨论因果矛盾的条件。
\Diagram{fig:m1antiq}{假设的超光速信号是类空路径；此图不表示真实光线}{\draw[axis](-.3,0)--(3.7,0)node[right]{$x/L$};\draw[axis](0,-1.2)--(0,3.4)node[above]{$ct/L$};\draw[curve](1,0)--(3.7,3);\draw[->,very thick](0,0)--(1.818182,.909091);\draw[->,very thick](1.818182,.909091)--(0,-.413223);\node[right]at(1.818182,.909091){$R$};\node[left]at(0,-.413223){$C$};}
\end{exercise}
\begin{solution}
出发信号$x=ut$与飞船相交，图\ref{fig:m1antiq}的接收事件为$t_R=L/(u-V)$、$x_R=uL/(u-V)$。以飞船系中从回复发出起经过的坐标时间$s>0$作参数，回复有$\Delta t'=s$、$\Delta x'=-us$。逆变换逐项给
\[
 \Delta t=\gamma(1-Vu/c^2)s,\qquad
 \Delta x=\gamma(V-u)s.
\]
回复到达地球要求$x_R+\Delta x=0$，所以$s=x_R/[\gamma(u-V)]>0$。代入时间增量，
\[
 t_C=t_R+\frac{(1-Vu/c^2)x_R}{u-V}
 =\frac{L[2u-V(1+u^2/c^2)]}{(u-V)^2}.
\]
当$Vu=c^2$，回复路径在实验室的时间增量为零，仍由以上参数式处理；直接以$\Delta x/\Delta t$表示斜率则会出现零分母。当$Vu>c^2$，沿回复路径的实验室时间增量为负，不能再用面向未来的普通速度图像解释其方向。
\Diagram{fig:m1antia}{若回复早于出发，地球可在发送之前收到由该发送触发的回信}{\draw[axis](0,-1.4)--(0,2)node[above]{地球钟};\fill(0,0)circle(1.5pt)node[right]{$A:t=0$发送};\fill(0,-.9)circle(1.5pt)node[right]{$C:t_C<0$收到回复};\draw[->,thick](-.5,-.8)--(-.5,-.1);}
因此$V>2u/(1+u^2/c^2)$时$t_C<0$。对$u>c$，这个门槛小于$c$，可由亚光速飞船达到。图\ref{fig:m1antiq}取$u=2c,V=.9c$作示意，回复位于地球出发事件的过去。图\ref{fig:m1antia}所示矛盾依赖“每个惯性系都能双向按同一超光速规则发信号”的假设；视超光速图样不具有这种信号能力，不能用它构造本装置。
\end{solution}

\begin{exercise}\label{ex:m1clocktrain}
列车固有长$L$，以速度$V$运动。车头和车尾各有一只在车系同步的钟。地面在$t=0$同时读取两钟，求车头减车尾的读数差，并求地面收到这两次读数的先后是否必然相同。
\Diagram{fig:m1trainq}{同时读取运动列车的两端钟}{\draw[axis](-.5,0)--(5,0)node[right]{$x$};\draw[curve](0,.5)rectangle(3.5,1.3);\node[above]at(0,1.3){尾};\node[above]at(3.5,1.3){头};\draw[axis](3.8,.9)--(4.8,.9)node[above]{$V$};\draw[guide](0,2)--(3.5,2);\node[above]at(1.75,2){$t=0$};}
\end{exercise}
\begin{solution}
图\ref{fig:m1trainq}中的同时截面长为$L/\gamma$。用$\Delta t'=\gamma(\Delta t-V\Delta x/c^2)$，得到头钟减尾钟
\[
 \Delta t'=-\gamma\frac{V}{c^2}\frac{L}{\gamma}=-\frac{VL}{c^2}.
\]
头钟在地面同时截面上落后。读数事件之间是类空间隔，其光信号到达同一观察者的时间还要加各自传播时间。
\Diagram{fig:m1traina}{地面接收者的位置影响读数信号的到达次序}{\draw[curve](0,0)--(4,0);\fill(0,0)circle(1.4pt)node[below]{尾};\fill(4,0)circle(1.4pt)node[below]{头};\fill(1,1.5)circle(1.4pt)node[above]{$D$};\draw[->](0,0)--(1,1.5);\draw[->](4,0)--(1,1.5);\node[below]at(2,-.5){$\Delta t_{\rm arr}=(r_{\text{头}}-r_{\text{尾}})/c$};}
图\ref{fig:m1traina}中的接收者若更靠近车尾，先收到尾部信号；靠近车头则相反；在垂直平分线上同时收到。钟面读数的差与信号到达次序没有一一对应关系。
\end{solution}

\begin{exercise}\label{ex:m1lifetime}
不稳定粒子的平均固有寿命为$\tau_0$，在实验室以$v$匀速飞行。求在距离$L$处仍未衰变的比例。再在粒子系说明为什么得到同一结果。
\Diagram{fig:m1lifeq}{粒子从产生点飞向距离为$L$的探测平面}{\draw[axis](-.3,0)--(5,0)node[right]{$x$};\fill(0,0)circle(2pt)node[below]{$O$};\draw[curve](4,-.7)--(4,1.2)node[above]{探测面};\draw[->,thick](.4,.3)--(2,.3)node[above]{$v$};\draw[<->](0,-.5)--(4,-.5)node[midway,below]{$L$};}
\end{exercise}
\begin{solution}
图\ref{fig:m1lifeq}中实验室飞行时间为$t=L/v$，粒子经历的固有时为$\tau=t/\gamma_v$。恒定衰变概率对应$\dd N/\dd\tau=-N/\tau_0$。分离变量并用$N(0)=N_0$，
\[
 \int_{N_0}^{N(\tau)}\frac{\dd n}{n}
 =-\frac1{\tau_0}\int_0^\tau\dd s,
 \qquad \ln\frac{N(\tau)}{N_0}=-\frac{\tau}{\tau_0}.
\]
取指数并代入飞行固有时，得
\[
 \frac N{N_0}=\exp\left[-\frac{L}{\gamma_vv\tau_0}\right].
\]
\Diagram{fig:m1lifea}{粒子系中探测面向静止粒子运动}{\fill(0,0)circle(2pt)node[above]{粒子};\draw[curve](3,-.7)--(3,1.2);\draw[axis](2.7,.5)--(1.5,.5)node[midway,above]{$v$};\draw[<->](0,-.5)--(3,-.5)node[midway,below]{$L/\gamma_v$};}
在图\ref{fig:m1lifea}的粒子系，产生点到探测面在同一时刻的距离缩短为$L/\gamma_v$，探测面以$v$靠近，所需时间$L/(\gamma_vv)$，恰为上面的固有时间。这里“平均寿命”不代表所有粒子都在同一时刻衰变，结果是统计存活比例。
\end{solution}

\begin{exercise}\label{ex:m1crossvelocity}
飞船$S'$以$V$相对实验室运动，在船上沿$y'$方向发射速度为$u$的粒子。求实验室速度、洛伦兹因子及粒子运动到船系坐标$y'=H$时的实验室位置。
\Diagram{fig:m1crossq}{船内的横向发射}{\draw[curve](0,0)rectangle(3,2.2);\fill(1,0)circle(1.5pt);\draw[axis](1,0)--(1,1.8)node[left]{$u$};\node[right]at(1,2){$H$};\draw[axis](3.4,1)--(4.6,1)node[right]{$V$};}
\end{exercise}
\begin{solution}
图\ref{fig:m1crossq}中$v'_x=0,v'_y=u$。逆速度变换给$v_x=V,v_y=u/\gamma$。于是
\[
 1-v^2/c^2=(1-\beta^2)(1-u^2/c^2),\quad \gamma_v=\gamma\gamma_u.
\]
到达上边界的事件在船系为$t'=H/u,x'=0,y'=H$。逆坐标变换得$t=\gamma H/u$，$x=\gamma VH/u$，$y=H$。
\Diagram{fig:m1crossa}{实验室中粒子轨迹斜向前方，与船保持相同纵向速度}{\draw[axis](-.3,0)--(4.8,0)node[right]{$x$};\draw[axis](0,-.2)--(0,2.5)node[above]{$y$};\draw[very thick,->](0,0)--(3.7,2);\draw[guide](0,2)--(3.7,2)--(3.7,0);\node[left]at(0,2){$H$};\node[below]at(3.7,0){$\gamma VH/u$};}
图\ref{fig:m1crossa}的斜率$H/(\gamma VH/u)=u/(\gamma V)$也等于$v_y/v_x$。若直接使用伽利略合成$v_y=u$，会使坐标变换与轨迹速度互相矛盾。
\end{solution}

\begin{exercise}\label{ex:m1coneenergy}
静止系各向同性的单次闪光总能量为$E_0$，各方向的光谱相同。实验室观察前向锥$\theta\leq\arccos\beta$。求锥内光子数与能量各占整次闪光的比例。
\Diagram{fig:m1coneq}{前向锥恰好由静止系前半球变换而来}{\draw[axis](-.3,0)--(4,0)node[right]{$x$};\draw[curve](0,0)--(3,1.3);\draw[curve](0,0)--(3,-1.3);\draw[guide](3,-1.3)--(3,1.3);\node[right]at(3,1.3){$\cos\theta_c=\beta$};}
\end{exercise}
\begin{solution}
由例\ref{ex:m1aberration}，图\ref{fig:m1coneq}中的光锥来自$0\leq\mu'\leq1$。静止系光子数各向同性，变换不改变光子数，所以数目占比为$1/2$。
\Diagram{fig:m1conea}{在静止系积分能量权重，权重为$\gamma(1+\beta\mu')$}{\draw[axis](0,0)--(3.8,0)node[right]{$\mu'$};\draw[axis](0,0)--(0,2.8)node[above]{能量权重};\draw[curve](0,1)--(3,2.3);\draw[guide](3,0)--(3,2.3);\node[below]at(3,0){$1$};\node[left]at(0,1){$\gamma$};}
能量按图\ref{fig:m1conea}的权重变换，锥内能量为
\[
 E_c=\frac{\gamma E_0}{2}\int_0^1(1+\beta\mu')\dd\mu'
 =\frac{\gamma E_0}{2}\left[\mu'+\frac\beta2\mu'^2\right]_0^1
 =\gamma E_0\left(\frac12+\frac\beta4\right).
\]
整个闪光的能量为$\gamma E_0$，故能量份额为$1/2+\beta/4$。即使高相对论速度，静止系前半球所含光子的实验室能量份额也只趋向$3/4$；这与小立体角内较大的角能量密度并不冲突。
\end{solution}

\begin{exercise}\label{ex:m1transverse}
运动光源的静止频率为$\nu_0$。分别求实验室观察方向与速度垂直、以及光源系发射方向与速度垂直时的实验室频率。
\Diagram{fig:m1transq}{“横向”可指两个不同惯性系中的方向}{\draw[axis](0,0)--(3.4,0)node[right]{$V$};\draw[axis](0,0)--(0,2)node[above]{$\theta=\pi/2$};\draw[->,thick](0,0)--(2.6,1.6)node[right]{$\theta'=\pi/2$的光};}
\end{exercise}
\begin{solution}
图\ref{fig:m1transq}中竖直的实验室光线有$\cos\theta=0$，所以$\nu=\nu_0/\gamma$，发生红移。若取源系$\cos\theta'=0$，则由逆能量变换$\nu=\gamma\nu_0$，发生蓝移。
\Diagram{fig:m1transa}{两个横向条件对应不同实验室角度}{\draw[axis](0,0)--(3.5,0)node[right]{$x$};\draw[curve](0,0)--(0,2);\draw[curve](0,0)--(2.6,1.6);\node[left]at(0,1.7){$\nu_0/\gamma$};\node[right]at(2.6,1.6){$\gamma\nu_0$};\draw[->](.9,0)arc(0:31.6:.9);\node[above]at(1.3,1.1){$\cos\theta=\beta$};}
第二种光线在实验室满足$\cos\theta=\beta$，如图\ref{fig:m1transa}。将其代入$\nu=\nu_0/[\gamma(1-\beta\cos\theta)]$，分母成为$\gamma(1-\beta^2)=1/\gamma$，与直接变换吻合。叙述横向多普勒效应时必须说明横向角在哪个系定义。
\end{solution}

\begin{exercise}\label{ex:m1fourwork}
静质量$m$的粒子由静止开始受恒定三力$f$沿$x$加速。求其动量、能量、速度和位移随实验室时间的变化，以及经历的固有时。
\Diagram{fig:m1workq}{实验室恒定三力的直线运动}{\draw[axis](-.3,0)--(4.7,0)node[right]{$x$};\fill(0,0)circle(2pt)node[below]{$m$};\draw[very thick,->](.2,.5)--(2,.5)node[midway,above]{$f$};}
\end{exercise}
\begin{solution}
以图\ref{fig:m1workq}的出发事件为$t=x=0$。恒三力满足$\dd p/\dd t=f$；用$p(0)=0$，有$\int_0^{p(t)}\dd p_1=f\int_0^t\dd t_1$，即$p=ft$。质壳关系和$v=pc^2/E$给
\[
 E=mc^2\sqrt{1+s^2},\qquad
 v=\frac{cs}{\sqrt{1+s^2}},\qquad s=\frac{ft}{mc}.
\]
位移的定积分为
\begin{align*}
 x(t)&=\int_0^t v(t_1)\dd t_1
 =\frac{mc^2}f\int_0^s\frac{q\dd q}{\sqrt{1+q^2}}\\
 &=\frac{mc^2}f[\sqrt{1+q^2}]_0^s
 =\frac{mc^2}f[\sqrt{1+s^2}-1].
\end{align*}
于是$fx=E-mc^2$，外力做功恰等于动能增加。
\Diagram{fig:m1worka}{恒力世界线渐近于光线，虚线斜率为1}{\draw[axis](-.2,0)--(4.5,0)node[right]{$fx/(mc^2)$};\draw[axis](0,-.2)--(0,4.5)node[above]{$ft/(mc)$};\draw[curve,domain=0:4,samples=80]plot({sqrt(1+\x*\x)-1},\x);\draw[guide](0,1)--(3.4,4.4);}
固有时间应沿图\ref{fig:m1worka}的加速世界线累计。由$\gamma_v=\sqrt{1+s^2}$，
\[
 \tau=\frac{mc}f\int_0^s\frac{\dd q}{\sqrt{1+q^2}}
 =\frac{mc}f\operatorname{arsinh}s.
\]
最后一步令$q=\sinh\eta$，则$\dd q=\cosh\eta\dd\eta$、$\sqrt{1+q^2}=\cosh\eta$，上下限为$0,\operatorname{arsinh}s$，积分为$[\eta]_0^{\operatorname{arsinh}s}$。反解$s=\sinh(f\tau/mc)$，代回得到
\[
 ct=\frac{mc^2}f\sinh\frac{f\tau}{mc},\qquad
 x=\frac{mc^2}f\left(\cosh\frac{f\tau}{mc}-1\right).
\]
从$v(t)=cs(1+s^2)^{-1/2}$用乘积与链式法则求导，
\[
 \frac{\dd v}{\dd t}=c\frac{\dd s}{\dd t}
 \left[(1+s^2)^{-1/2}-s^2(1+s^2)^{-3/2}\right]
 =\frac fm(1+s^2)^{-3/2}.
\]
纵向固有加速度为$\gamma_v^3\dd v/\dd t=f/m$。恒三力对应恒纵向固有加速度，实验室加速度则不断减小。$s\ll1$时$x\simeq ft^2/(2m)$、$\tau\simeq t$，任何有限$t$都有$v<c$。
\end{solution}
\begin{exercise}\label{ex:m1twoemitters}
静止系中两灯位于$\pm L/2$，同时向中点发光。求两束光相遇事件在运动系的坐标，以及两发射事件在运动系是否同时。
\Diagram{fig:m1emitq}{两灯同时发射相向光线}{\draw[axis](-2.3,0)--(2.5,0);\fill(-2,0)circle(1.5pt)node[below]{$A$};\fill(2,0)circle(1.5pt)node[below]{$B$};\draw[->](-1.8,.5)--(-.2,.5);\draw[->](1.8,.5)--(.2,.5);}
\end{exercise}
\begin{solution}
在图\ref{fig:m1emitq}的$S$系中，左右发射事件分别为$A:(0,-L/2)$、$B:(0,L/2)$，坐标次序为$(ct,x)$。两束光到中点均走$L/2$，故相遇于$C:(L/2,0)$。变换至沿$+x$以$V$运动的$S'$，得到
\begin{align*}
 A'&:(\gamma\beta L/2,-\gamma L/2),\\
 B'&:(-\gamma\beta L/2,\gamma L/2),\\
 C'&:(\gamma L/2,-\gamma\beta L/2).
\end{align*}
所以右灯先发光，$t'_B-t'_A=-\gamma\beta L/c$。
\Diagram{fig:m1emita}{运动系中先发射的光走得更久，相遇点仍唯一}{\draw[axis](-2.7,0)--(2.5,0)node[right]{$x'$};\draw[axis](0,-1.7)--(0,3)node[above]{$ct'$};\draw[very thick](-2,1.2)--(-1.2,2)--(2,-1.2);\node[left]at(-2,1.2){$A$};\node[right]at(2,-1.2){$B$};\node[above]at(-1.2,2){$C$};}
图\ref{fig:m1emita}中，左行和右行光仍具有光速。由$A'$到$C'$，
\[
 c(t'_C-t'_A)=\frac{\gamma L}2(1-\beta)
 =x'_C-x'_A;
\]
由$B'$到$C'$，
\[
 c(t'_C-t'_B)=\frac{\gamma L}2(1+\beta)
 =x'_B-x'_C.
\]
两条光路的传播时间不同，正好补偿了两次发射的不同步。相遇是同一时空事件，不会因参考系改变而消失。$x'_B-x'_A=\gamma L$是两个不同时发射事件的位置差；两灯在$S'$的同时距离仍为$L/\gamma$。
\end{solution}

\begin{exercise}\label{ex:m1clockoffset}
两个静止于$S'$的钟相距$L$，钟面相对坐标时间偏移分别为$b_A,b_B$。实验室在$t=0$同时读取，求读数差；什么偏移能使这次读数相同？
\Diagram{fig:m1offsetq}{钟面读数含可人为设置的常数偏移}{\draw[curve](0,0)circle(.5);\draw[curve](4,0)circle(.5);\node at(0,0){$A$};\node at(4,0){$B$};\node[above]at(0,.7){$t'+b_A$};\node[above]at(4,.7){$t'+b_B$};\draw[axis](1,-1)--(2.5,-1)node[right]{$V$};}
\end{exercise}
\begin{solution}
在图\ref{fig:m1offsetq}的实验室同时截面，两钟坐标时间差$\Delta t'=-VL/c^2$，因此钟面差$T_B-T_A=b_B-b_A-VL/c^2$。使本次读数相同需$b_B-b_A=VL/c^2$。
\Diagram{fig:m1offseta}{两钟速率相同，偏移只改变读数的常数差}{\draw[axis](0,0)--(4.5,0)node[right]{$t$};\draw[axis](0,0)--(0,3.6)node[above]{读数};\draw[curve](0,.4)--(3.8,2.8);\draw[guide](0,1)--(3.8,3.4);\node[right]at(3.8,2.8){固定偏移};}
图\ref{fig:m1offseta}中两钟在实验室的速率均为$1/\gamma$，所以在每个实验室同时截面上，该读数差都相同；这个偏移相当于按实验室同时性而非车系爱因斯坦同时性设置钟。钟的正常走时速率与一组钟的同步约定是不同问题。
\end{solution}

\begin{exercise}\label{ex:m1propersegments}
粒子先以$v_1$运动实验室时间$T_1$，再以$v_2$运动$T_2$，方向任意，变速时间忽略。求总固有时；若已知两端点坐标差，何时可以用端点间隔代替？
\Diagram{fig:m1segmentsq}{折线世界线的两段固有时须相加}{\STAxes\draw[very thick](0,0)--(2,2.5)--(1,4);\draw[guide](0,0)--(1,4);\node[right]at(2,2.5){$B$};\node[right]at(1,4){$C$};}
\end{exercise}
\begin{solution}
图\ref{fig:m1segmentsq}每段匀速，故$\tau=T_1\sqrt{1-v_1^2/c^2}+T_2\sqrt{1-v_2^2/c^2}$。端点间隔给的是直接连接两端的匀速世界线固有时$\tau_{AC}=\sqrt{(T_1+T_2)^2-|\vct v_1T_1+\vct v_2T_2|^2/c^2}$。
\Diagram{fig:m1segmentsa}{在端点共同静止系中，直线停留的钟积累时间最多}{\draw[axis](-1.5,0)--(2,0)node[right]{$x_*$};\draw[axis](0,0)--(0,4.3)node[above]{$ct_*$};\draw[very thick](0,0)--(1.3,2.1)--(0,3.8);\draw[guide](0,0)--(0,3.8);}
在图\ref{fig:m1segmentsa}的端点静止系，实际轨迹$\int\sqrt{1-v_*^2/c^2}\dd t_*\leq\int\dd t_*=\tau_{AC}$，等号要求全程$v_*=0$，即两段本来就是同一匀速速度矢量。仅两段速率相同而方向改变，仍不能用端点间隔代替曲线积分。
\end{solution}

\begin{exercise}\label{ex:m1inversevel}
实验室测得粒子$v_x=.8c,v_y=.3c$。求使其纵向速度为零的$x$向惯性系，以及该系中的横向速度和洛伦兹因子。
\Diagram{fig:m1inverseq}{选择变换速度消去纵向分量}{\draw[axis](0,0)--(4,0)node[right]{$x$};\draw[axis](0,0)--(0,2.3)node[above]{$y$};\draw[->,thick](0,0)--(3.2,1.2);\draw[guide](3.2,0)--(3.2,1.2);}
\end{exercise}
\begin{solution}
图\ref{fig:m1inverseq}中，设所求参考系沿$x$以$V$运动。速度变换的分母$1-Vv_x/c^2>0$，所以$v'_x=0$等价于$V=v_x=.8c$，其洛伦兹因子为$\gamma=5/3$。横向速度为
\[
 v'_y=\frac{v_y}{\gamma(1-Vv_x/c^2)}
 =\frac{.3c}{(5/3)(1-.64)}=.5c.
\]
因此新系中粒子仅沿$y'$运动，$\gamma_{v'}=(1-.25)^{-1/2}=2/\sqrt3$。
\Diagram{fig:m1inversea}{新系中轨迹纯横向，速率仍小于$c$}{\draw[axis](0,0)--(3.7,0)node[right]{$x'$};\draw[axis](0,0)--(0,2.7)node[above]{$y'$};\draw[very thick,->](0,0)--(0,2.1)node[right]{$.5c$};}
横向速度增大的原因可由微小事件差看出：$\dd y'=\dd y$，而
\[
 \dd t'=\gamma(1-Vv_x/c^2)\dd t=.6\dd t.
\]
相同横向位移除以更短的新系坐标时间，得到图\ref{fig:m1inversea}的$.5c$，并非横向位移本身被放大。四速度提供独立核对：原速率平方$.73c^2$，故
\[
 \gamma_v=\frac1{\sqrt{.27}},\qquad
 \gamma_{v'}=\gamma\gamma_v(1-.64)
 =\frac{.6}{\sqrt{.27}}=\frac2{\sqrt3}.
\]
只消去纵向速度仍未使粒子静止。若再在$S'$中沿$y'$以$.5c$变换，才到达粒子的静止系；两次非共线变换的空间轴关系须另行确定。
\end{solution}

\begin{exercise}\label{ex:m1numberangular}
各向同性光源在固有时内每秒发射$\dot N_0$个光子。求实验室按发射时间定义的角光子率，积分求总光子率，并与总功率不变比较。
\Diagram{fig:m1numberq}{同一批光子的数目不变，角密度与时间率改变}{\draw[axis](0,0)--(3.8,0)node[right]{$V$};\draw[curve](0,0)--(3,1);\draw[curve](0,0)--(3,1.5);\node[right]at(3,1.5){$\dd\Omega$};}
\end{exercise}
\begin{solution}
图\ref{fig:m1numberq}中，源系在$\dd\tau$内向$\dd\Omega'$发出的光子数为$\dd N=\dot N_0\dd\tau\dd\Omega'/(4\pi)$。令$\delta=[\gamma(1-\beta\cos\theta)]^{-1}$。由光行差，
\[
 \frac{\dd\cos\theta}{\dd\cos\theta'}
 =\frac1{\gamma^2(1+\beta\cos\theta')^2}=\delta^{-2},
\]
方位角不变，所以$\dd\Omega=\delta^{-2}\dd\Omega'$。再用实验室发射时间$\dd t=\gamma\dd\tau$，得到
\[
 \frac{\dd\dot N}{\dd\Omega}=\frac{\dot N_0}{4\pi\gamma}\delta^2.
\]
\Diagram{fig:m1numbera}{单次光子数不变，时间膨胀降低实验室总发射率}{\draw[axis](0,0)--(4.2,0)node[right]{$\beta$};\draw[axis](0,0)--(0,2.7)node[above]{$\dot N/\dot N_0$};\draw[curve,domain=0:1,samples=70]plot({3.5*\x},{2.3*sqrt(1-\x*\x)});\node[left]at(0,2.3){$1$};}
全角积分可以直接计算。以$\mu=\cos\theta$，
\begin{align*}
 \int\delta^2\dd\Omega
 &=\frac{2\pi}{\gamma^2}\int_{-1}^1\frac{\dd\mu}{(1-\beta\mu)^2}\\
 &=\frac{2\pi}{\gamma^2}
 \left[\frac{1}{\beta(1-\beta\mu)}\right]_{-1}^1\\
 &=\frac{2\pi}{\beta\gamma^2}
 \left[\frac1{1-\beta}-\frac1{1+\beta}\right]=4\pi.
\end{align*}
在$\beta=0$处取连续极限。因此图\ref{fig:m1numbera}给出总光子率$\dot N=\dot N_0/\gamma$。若源系每光子能量为$E_0$，同一批各向同性光子的总动量为零，整体能量变换便给平均能量$\langle E\rangle=\gamma E_0$。于是$P=\dot N\langle E\rangle=\dot N_0E_0$。按实验室发射时间计，光子数率因时间膨胀减小，平均单光子能量增加，两者相抵。接收时间率还须计入传播延迟。
\end{solution}

\begin{exercise}\label{ex:m1redblue}
单色运动源静止频率$\nu_0$。求实验室观察到无频移的方向，并确定蓝移、红移的角域。
\Diagram{fig:m1redq}{多普勒因子随实验室观察角变化}{\draw[axis](0,0)--(4,0)node[right]{$V$};\draw[curve](0,0)--(2.3,1.6);\draw[curve](0,0)--(-1,1.6);\draw[guide](0,0)--(0,2.2);\node at(1,.35){$\theta$};}
\end{exercise}
\begin{solution}
图\ref{fig:m1redq}中，$0<\beta<1$时$\nu/\nu_0=[\gamma(1-\beta\cos\theta)]^{-1}$。令此比为1，得到
\[
 \cos\theta_* =\frac{1-1/\gamma}{\beta}
 =\frac{\gamma-1}{\gamma\beta}.
\]
利用$(\gamma-1)(\gamma+1)=\gamma^2\beta^2$，可改写为$\cos\theta_*=\gamma\beta/(\gamma+1)$。该值介于0和1之间，对应唯一的前向无频移圆锥。
\Diagram{fig:m1reda}{无频移锥内蓝移，锥外红移}{\draw[axis](-1,0)--(4,0)node[right]{$x$};\draw[curve](0,0)--(1.8,2);\draw[curve](0,0)--(1.8,-2);\node[right]at(2,.5){蓝移};\node[left]at(-.2,1.5){红移};\node[above]at(1.8,2){$\theta_*$};}
图\ref{fig:m1reda}的锥内满足$\cos\theta>\cos\theta_*$，多普勒分母小于1，因而蓝移；锥外红移。将边界方向反变换到源系，有
\[
 \cos\theta'_* =\frac{\cos\theta_*-\beta}{1-\beta\cos\theta_*}
 =-\frac{\gamma\beta}{\gamma+1}=-\cos\theta_*.
\]
源系略向后的发射方向也可能在实验室无频移。低速展开$\gamma=1+\beta^2/2+\cdots$给$\cos\theta_*\simeq\beta/2$，边界略小于$\pi/2$；高速时$\cos\theta_*\to1$。当$\beta=0$时所有方向都无频移，不能把$\beta\to0$所得的边界锥误认为静止光源只有一个无频移方向。
\end{solution}


